\section{Introduction}

Is DeepSeek-R1 39\% faithful or 95\%? Chen et al.~\cite{chen2025reasoning} report that DeepSeek-R1 acknowledges sycophantic hints in its chain-of-thought only 39\% of the time. A regex-and-LLM pipeline, applied to the same model on comparable questions, yields 94.8\%. An independent Claude Sonnet~4 judge returns 74.8\%. Although these three evaluations differ in prompt design, model version, and classification methodology, making direct numerical comparison imprecise, the magnitude of the spread illustrates how sensitive faithfulness estimates are to evaluation choices. Which number is right?

Chain-of-thought prompting~\cite{wei2022cot} has become the dominant paradigm for eliciting reasoning from large language models, and a growing literature treats measured faithfulness rates as objective properties of models~\cite{turpin2023unfaithful, lanham2023measuring, chen2025reasoning}. Papers report single numbers and compare them across models, drawing conclusions about which architectures or training methods produce more faithful reasoning~\cite{feng2025more, arcuschin2025wild}. Yet these numbers are products of two things: the model's actual behavior \emph{and} the classifier used to judge that behavior. When classifiers disagree substantially, faithfulness comparisons become classifier comparisons in disguise.

This problem has well-established precedents in the measurement sciences. Jacovi and Goldberg~\cite{jacovi2020faithfulness} argued that faithfulness itself admits multiple valid definitions, and Parcalabescu and Frank~\cite{parcalabescu2023measuring} showed that different operationalizations yield divergent measurements. More broadly, Gordon et al.~\cite{gordon2021disagreement} demonstrated that correcting for annotator disagreement in toxicity classification reduced apparent performance from 0.95 to 0.73 ROC AUC, Plank~\cite{plank2022problem} argued that such label variation is an inherent property of subjective tasks rather than noise, and Bean et al.~\cite{bean2025measuring} found systematic construct validity failures across 445 LLM benchmarks. The use of LLMs as automated judges~\cite{zheng2023judging} introduces additional classifier-specific biases that have been catalogued by Gu et al.~\cite{gu2024survey} and Ye et al.~\cite{ye2024justice}. Despite these known concerns, CoT faithfulness studies routinely report single-classifier results without sensitivity analysis.

This paper provides an initial quantification of the problem. Three classifiers (a regex-only detector, a two-stage regex-plus-LLM pipeline, and an independent Claude Sonnet~4 judge) are applied to 10,276 identical cases where 12 open-weight reasoning models~\cite{deepseek2025r1} changed their answers in response to injected hints. While the underlying data are drawn from a companion empirical study~\cite{young2026lietome}, the present paper addresses a distinct research question: not \emph{how faithful} are these models, but \emph{how stable} are faithfulness measurements across classification methodologies. This question requires a different analytical framework (inter-classifier agreement, construct comparison) than the cross-family analysis of the parent study. Four results challenge the practice of reporting single faithfulness numbers:

\begin{enumerate}
    \item \textbf{Large aggregate disagreement.} Overall faithfulness rates range from 74.4\% (regex-only) through 82.6\% (pipeline) to 69.7\% (Sonnet judge). Per-model gaps range from 2.6 to 30.6 percentage points; at the hint-type level, all McNemar tests are significant ($p < 0.001$).
    \item \textbf{Hint-type-specific disagreement.} The gap is not uniform. Sycophancy hints~\cite{sharma2023sycophancy} show a 43.4-percentage-point classifier gap ($\kappa = 0.06$, ``slight'' agreement), while grader hints show only 2.9 percentage points ($\kappa = 0.42$, ``moderate'' agreement). Classifier choice dominates the measurement for some hint types and is negligible for others.
    \item \textbf{Asymmetric confusion structure.} The disagreements are directional, not random. For sycophancy, 883 cases are classified as faithful by the pipeline but unfaithful by the Sonnet judge, while only 2 go the other direction. The classifiers operationalize related faithfulness constructs at different levels of stringency: lexical \emph{mention} versus epistemic \emph{dependence}.
    \item \textbf{Ranking reversals.} Qwen3.5-27B~\cite{qwen2024technical} ranks 1st by the pipeline (98.9\% faithful) but 7th by Sonnet (68.3\%); OLMo-3.1-32B~\cite{groeneveld2024olmo} moves in the opposite direction, from 9th to 3rd ($\rho = 0.67$, $p = 0.017$). Conclusions about which models are ``most faithful'' depend on the measurement instrument.
\end{enumerate}

Three hypotheses are tested:
\begin{description}
\item[H1:] Aggregate faithfulness rates differ significantly across the three classifiers when applied to the same data, reflecting construct-level differences rather than random measurement error.
\item[H2:] The magnitude of classifier disagreement varies by hint type, with social-pressure hints (sycophancy, consistency) producing larger gaps than rule-based hints (grader, unethical).
\item[H3:] Classifier choice can reverse the ordinal ranking of models, such that conclusions about which models are ``most faithful'' depend on the measurement instrument.
\end{description}

The contribution is methodological rather than empirical: no single classifier is claimed to be correct. Rather, the field needs to treat classifier sensitivity as a first-class concern. Just as Gordon et al.~\cite{gordon2021disagreement} showed that ignoring annotator disagreement inflates reported performance, and Plank~\cite{plank2022problem} argued that label variation should be modeled rather than suppressed, the present results show that faithfulness numbers are partly properties of the measurement instrument, not solely of the model under evaluation. Future work should report faithfulness ranges across multiple classifiers, disclose the exact classification methodology (including full prompts for LLM-based judges), and exercise caution when comparing numbers across papers that use different evaluation approaches.
